Every convention below is pinned as a named constant in
\texttt{spinor\_trace\_bridge/src/sdbridge/conventions.py}. Downstream code
imports the constant and may not re-derive a sign locally; a change to that file
is a scientific change and triggers a full re-run.

\paragraph{Metric.} Signature $(-,+,\dots,+)$, index $0$ timelike. The tensor
implementation's own metric function returns exactly this, and a test asserts the
equality rather than trusting it.

\paragraph{Epsilon and Hodge star.} $\epsilon_{01\ldots 9}=+1$;
$(\star F)_I = \frac{1}{(d-p)!}\epsilon_I{}^J F_J$, implemented on sorted
multi-indices as $\pm$ the complementary component times the raising signs, so
that no $10^{10}$ epsilon tensor is ever built. The value of $\star^2$ is
\emph{computed} at runtime, not assumed, and asserted to be $+\starSquared$.

\paragraph{Self-duality.} $\star F = +F$; projector $(1+\star)/2$, with $1/2$
taken as the modular inverse of $2$. Self-dual and anti-self-dual subspaces are
each $\dimSelfdual$-dimensional.

\paragraph{Clifford algebra.} $\{\Gamma^\mu,\Gamma^\nu\}=2\eta^{\mu\nu}$.
Chirality: $\Lambda^{\mathrm{even}}W$ is the positive-chirality module. The
Chevalley pairing carries the reversal sign $(-1)^{p(p-1)/2}$, which is what
makes $\sigma^\mu_{ab}$ symmetric; the symmetry is asserted, not assumed.

\paragraph{Antisymmetrised five-gamma.}
$(\Gamma^{\mu_1\ldots\mu_5})_{ab} = \frac{1}{5!}\sum_{\pi}\mathrm{sgn}(\pi)\,
(\sigma\bar\sigma\sigma\bar\sigma\sigma)_{ab}$, unit weight on the identity
permutation.

\paragraph{Forward normalisation.} Equation~\eqref{eq:bridge} sums over all index
tuples with the $1/5!$; because both objects are totally antisymmetric this
equals the sum over sorted tuples with unit weight, which is what the code does.

\paragraph{Inverse.} No closed form is claimed. A left inverse on the self-dual
subspace is constructed by inverting the $\dimSelfdual\times\dimSelfdual$
restriction in an explicit pivot basis; its defining property
$\mathrm{inverse}\circ\mathrm{forward}=P_{\text{self-dual}}$ is checked exactly.

\paragraph{Coefficient field.} $\mathbb{F}_p$ for $p\in\{\bridgePrimes\}$. The
bridge contains no floating-point step and no tolerance parameter.

\paragraph{Spinor index placement.} Determined, not chosen: see
section~\ref{sec:bridge}. Recorded as an open confirmation item.


\paragraph{Formulation independence of the leading-degree statement.}
Section~\ref{sec:czero}'s closure bookkeeping uses the fact that the free
stress-tensor trace vanishes at quadratic order, so that $\operatorname{Tr}\tau$
first contributes at field degree four. The vanishing does not depend on the
improvement convention: for any
$T_{\mu\nu} \sim F_{\mu\ldots}F_\nu{}^{\ldots} - c\,\eta_{\mu\nu}\langle F,F\rangle$
the trace is $(1-cd)\langle F,F\rangle$, a multiple of the single quadratic
scalar $\langle F,F\rangle$; and $\langle F,F\rangle$ vanishes identically on
either eigenspace of the star, because $F\wedge F$ is a top form built from two
copies of an odd-degree form and $\langle F,\star F\rangle$ is proportional to
it. A covariant action for a chiral form requires auxiliary
structure~\cite{Pasti:1995tn,Pasti:1997gx}, but the argument above uses only
the quadratic scalar available to the free theory, so it is insensitive to which
such formulation is adopted. Confirmation that the intended conventions are
those used here is coauthor review item~G-10.
